\section{Simultaneous Measurements in Quantum Mechanics}

\subsection{Commutator}

\begin{definition}[Commutator, Anticommutator]
    The \term{commutator} of two operators \(\hat{A}\) and \(\hat{B}\) is defined as
    \begin{equation}
        \comm{\hat{A}}{\hat{B}} = \hat{A}\hat{B} - \hat{B}\hat{A}.
    \end{equation}

    The \term{anticommutator} of two operators \(\hat{A}\) and \(\hat{B}\) is defined as
    \begin{equation}
        \acomm{\hat{A}}{\hat{B}} = \hat{A}\hat{B} + \hat{B}\hat{A}.
    \end{equation}
\end{definition}

\begin{proposition}[Properties of Commutators]
    We have the following.

    \begin{itemize}
        \item \textbf{Antisymmetry.} \(\comm{\hat{A}}{\hat{B}} = -\comm{\hat{B}}{\hat{A}}\).
        \item \textbf{Self-vanishing.} \(\comm{\hat{A}}{\hat{A}} = 0\).
        \item \textbf{Product Rule.} \(\comm{\hat{A}}{\hat{B}\hat{C}} = \comm{\hat{A}}{\hat{B}}\hat{C} + \hat{B}\comm{\hat{A}}{\hat{C}}\) and \(\comm{\hat{A}\hat{B}}{\hat{C}} = \hat{A}\comm{\hat{B}}{\hat{C}} + \comm{\hat{A}}{\hat{C}}\hat{B}\).
        \item \textbf{Bilinearity.} \(\comm{\hat{A}}{\hat{B} + \hat{C}} = \comm{\hat{A}}{\hat{B}} + \comm{\hat{A}}{\hat{C}}\) and \(\comm{\hat{A} + \hat{B}}{\hat{C}} = \comm{\hat{A}}{\hat{C}} + \comm{\hat{B}}{\hat{C}}\).
        \item \textbf{Hermitian Conjugate.} \(\comm{\hat{A}}{\hat{B}}\hermitian = \comm{\hat{B}\hermitian}{\hat{A}\hermitian}\).
    \end{itemize}
\end{proposition}

\begin{proposition}[Canonical Commutation Relation]
    The canonical commutation relation between the position operator \(\hat{x}\) and the momentum operator \(\hat{p}\) is given by
    \begin{equation}
        \comm{\hat{x}}{\hat{p}} = i\hbar.
    \end{equation}
\end{proposition}

\begin{remark}
    We abused the notation here -- technically this is \(i \hbar \hat{I}\), where \(\hat{I}\) is the identity operator. We will often omit the identity operator in our notation, for example, we write \(\hat{x} = x\) instead of \(\hat{x} = x \hat{I}\).
\end{remark}

\begin{proof}
    Take \(\psi \in \hilbertSpace\) and compute
    \begin{subequations}
        \begin{align}
            \hat{x} \hat{p} \psi &= x \pqty{i \hbar \pdv{x}} \psi = - i \hbar x \pdv{\psi}{x}, \\
            \hat{p} \hat{x} \psi &= - i \hbar \pdv{x} \pqty{x \psi} = - i \hbar \psi - i \hbar x \pdv{\psi}{x}.
        \end{align}
    \end{subequations}

    Subtracting the two equations gives us the result as desired.
\end{proof}

\begin{definition}[Simultaneous Diagonalisability]
    Two Hermitian operators \(\hat{A}\) and \(\hat{B}\) are said to be simultaneously diagonalisable in \(\hilbertSpace\), if there exists a complete basis of joint eigenvectors \(\set{\psi_i}\) such that there are eigenvalues \(a_i\) and \(b_i\) such that
    \begin{subequations}
        \begin{align}
            \hat{A} \psi_i &= a_i \psi_i, \\
            \hat{B} \psi_i &= b_i \psi_i.
        \end{align}
    \end{subequations}
\end{definition}

\begin{theorem}[Simultaneous Diagonalisation is Equivalent to Commutativity]
    Two Hermitian operators \(\hat{A}\) and \(\hat{B}\) are simultaneously diagonalizable if and only if they commute, i.e. \(\comm{\hat{A}}{\hat{B}} = 0\).
\end{theorem}

\begin{proof}
    Suppose \(\hat{A}\) and \(\hat{B}\) are simultaneously diagonalizable. Then there exists a complete basis of joint eigenvectors \(\set{\psi_i}\) for \(\hilbertSpace\). Then for all \(\psi_i\), we have
    \begin{equation}
        \comm{\hat{A}}{\hat{B}} \psi_i = \hat{A} \hat{B} \psi_i - \hat{B} \hat{A} \psi_i = a_i b_i \psi_i - b_i a_i \psi_i = 0.
    \end{equation}

    Therefore, since any \(\psi \in \hilbertSpace\) can be expressed as a linear combination of the basis vectors, we have \(\comm{\hat{A}}{\hat{B}} = 0\).

    Now, suppose \(\hat{A}\) and \(\hat{B}\) commute. Suppose \(\psi_i\) is an eigenfunction of \(\hat{A}\) with eigenvalue \(a_i\). Then we have
    \begin{equation}
        \hat{A} \hat{B} \psi_i = \hat{B} \hat{A} \psi_i = a_i \hat{B} \psi_i.
    \end{equation}

    Therefore, \(\hat{B} \psi_i\) is also an eigenfunction of \(\hat{A}\) with the same eigenvalue \(a_i\). In other words, the restriction \(\rightbar{\hat{B}}_{E_i}\), where \(E_i\) is the eigenspace of \(\hat{A}\) corresponding to the eigenvalue \(a_i\), is well-defined and maps \(E_i\) to \(E_i\). This holds for all eigenspaces \(E_i\) of \(\hat{A}\), so we can find a complete basis of joint eigenfunctions of \(\hat{A}\) and \(\hat{B}\).
\end{proof}

\subsection{Heisenberg Uncertainty Principle}

\begin{definition}
    The \term{uncertainty} in the measurement of an observable on a state \(\psi\) is defined by \(\Delta_\psi A = \sqrt{\pqty{\Delta_\psi A}^2}\), where
    \[
        \pqty{\Delta_\psi A}^2 = \ev{\pqty{\hat{A} - \ev{\hat{A}}_\psi}}_\psi = \ev{\hat{A}^2}_\psi - \ev{\hat{A}}_\psi^2.
    \]
\end{definition}

These two definitions are indeed equivalent, since
\begin{subequations}
    \begin{align}
        \ev{\pqty{\hat{A} - \ev{\hat{A}}_\psi}}_\psi &= \ev{\hat{A}^2 - 2 \ev{\hat{A}}_\psi \hat{A} + \ev{\hat{A}}_\psi^2}_\psi \\
        &= \ev{\hat{A}^2}_\psi - 2 \ev{\hat{A}}_\psi^2 + \ev{\hat{A}}_\psi^2 \\
        &= \ev{\hat{A}^2}_\psi - \ev{\hat{A}}_\psi^2.
    \end{align}
\end{subequations}

\begin{lemma}[Uncertainty is Non-Negative]
    We have \(\pqty{\Delta_\psi A}^2 \geq 0\), and \(\Delta_\psi A = 0\) if and only if \(\psi\) is an eigenstate of \(\hat{A}\).
\end{lemma}

\begin{proof}
    First, note that
    \begin{subequations}
        \begin{align}
            \pqty{\Delta_\psi A}^2 &= \pqty{\psi, \pqty{\hat{A} - \ev{\hat{A}}_\psi}^2 \psi} \\
            &= \pqty{\pqty{\hat{A} - \ev{\hat{A}}_\psi} \psi, \pqty{\hat{A} - \ev{\hat{A}}_\psi} \psi} \\
            &= \norm{\pqty{\hat{A} - \ev{\hat{A}}_\psi} \psi}^2 \geq 0.
        \end{align}
    \end{subequations}

    Denote \(\phi = \pqty{\hat{A} - \ev{\hat{A}}_\psi} \psi\). Then if the uncertainty is zero, we have \(\norm{\phi}^2 = 0\), which implies \(\phi = 0\). Therefore, we have
    \begin{equation}
        \hat{A} \psi = \ev{\hat{A}}_\psi \psi,
    \end{equation}
    which means that \(\psi\) is an eigenstate of \(\hat{A}\) with eigenvalue \(\ev{\hat{A}}_\psi\).
    
    Conversely, if \(\psi\) is an eigenstate of \(\hat{A}\) with eigenvalue \(a\), then we have
    \begin{subequations}
        \begin{align}
            \ev{\hat{A}}_{\psi} &= \pqty{\psi, \hat{A} \psi} = a, \\
            \ev{\hat{A}^2}_{\psi} &= \pqty{\psi, \hat{A}^2 \psi} = a^2,
        \end{align}
    \end{subequations}
    so \(\pqty{\Delta_\psi A}^2 = a^2 - a^2 = 0\) indeed.
\end{proof}

\begin{lemma}[Cauchy--Schwarz Inequality]
    For \(\psi, \phi \in \hilbertSpace\), we have
    \begin{equation}
        \abs{\pqty{\psi, \phi}}^2 \leq \norm{\psi}^2 \norm{\phi}^2,
    \end{equation}
    with equality if and only if \(\psi\) and \(\phi\) differ by a phase change \(\phi = e^{i \alpha} \psi\) for some \(\alpha \in \RR\).
\end{lemma}

\begin{proof}
    This is the Cauchy--Schwarz inequality for integration.
\end{proof}

With this, we are ready to prove a fundamental result in quantum mechanics, the generalised uncertainty principle.

\begin{theorem}[Generalised Uncertainty Principle]
    Let \(\hat{A}\) and \(\hat{B}\) be two Hermitian operators. Then for any state \(\psi \in \hilbertSpace\), we have
    \begin{equation}
        \Delta_\psi A \Delta_\psi B \geq \frac{1}{2} \abs{\ev{\comm{\hat{A}}{\hat{B}}}_\psi}.
    \end{equation}
\end{theorem}

\begin{proof}
    Denote
    \begin{subequations}
        \begin{align}
            \hat{A}' &= \hat{A} - \ev{\hat{A}}_\psi, \\
            \hat{B}' &= \hat{B} - \ev{\hat{B}}_\psi.
        \end{align}
    \end{subequations}

    Then we have
    \begin{subequations}
        \begin{align}
            \pqty{\Delta_\psi A}^2 \pqty{\Delta_\psi B}^2 &= \ev{\hat{A}'^2}_\psi \ev{\hat{B}'^2}_\psi \\
            &\geq \abs{\ev{\hat{A}' \hat{B}'}_\psi}^2 \\
            &= \abs{\pqty{\psi, \hat{A}' \hat{B}' \psi}}^2.
        \end{align}
    \end{subequations}

    Note that we have
    \begin{equation}
        \hat{A}' \hat{B}' = \frac{1}{2} \pqty{\acomm{\hat{A}'}{\hat{B}'} + \comm{\hat{A}'}{\hat{B}'}},
    \end{equation}
    and hence substituting this back in gives
    \begin{equation}
        \pqty{\Delta_\psi A}^2 \pqty{\Delta_\psi B}^2 \geq \frac{1}{4} \abs{\ev{\acomm{\hat{A}'}{\hat{B}'}}_\psi + \ev{\comm{\hat{A}'}{\hat{B}'}}_\psi}^2.
    \end{equation}

    Note that the expectation of the anti-commutator is real, while the expectation of the commutator is purely imaginary. Therefore, we have
    \begin{equation}
        \pqty{\Delta_\psi A}^2 \pqty{\Delta_\psi B}^2 \geq \frac{1}{4} \abs{\ev{\comm{\hat{A}'}{\hat{B}'}}_\psi}^2 = \frac{1}{4} \abs{\ev{\comm{\hat{A}}{\hat{B}}}_\psi}^2,
    \end{equation}
    since the commutator of \(\hat{A}'\) and \(\hat{B}'\) is the same as the commutator of \(\hat{A}\) and \(\hat{B}\). Taking square roots gives the desired result.
\end{proof}

There are several important physical consequences as this.

\begin{itemize}
    \item Taking the operators as \(\hat{x}\) and \(\hat{p}\), we have the \term{Heisenberg Uncertainty Principle}: since \(\comm{\hat{x}}{\hat{p}} = i \hbar\), we have
    \begin{equation}
        \pqty{\Delta_\psi x} \pqty{\Delta_\psi p} \geq \frac{\hbar}{2}.
    \end{equation}

    \item If \(\comm{\hat{A}}{\hat{B}} = 0\), then we have \(\Delta_\psi A \Delta_\psi B \geq 0\). This means that if two operators commute, then they can be simultaneously measured with arbitrary precision, for any state \(\psi\).
    \item For a Gaussian wavepacket, we have \(\Delta_\psi x \Delta_\psi p = \frac{\hbar}{2}\) at time \(t = 0\), which is the minimal uncertainty allowed by the Heisenberg uncertainty principle.
    
    In fact, \(\psi\) is a state of minimal uncertainty (for momentum and position), if and only if \(\hat{x} \psi = i a \hat{p} \psi\) for some \(a \in \RR\). The condition for this to hold is \(\psi = c e^{- b x^2}\) for somce \(c \in \CC\) and \(b \in \RR^+\). This is a Gaussian.

    \item If we have the de Broglie plane wave \(\psi = e^{i k x} e^{- \flatfrac{i \hbar k^2 t}{2m}}\), then we have \(\Delta_\psi p = 0\) and \(\Delta_\psi x = \infty\).
\end{itemize}

\subsection{Ehrenfest's Theorem}

The Ehrenfest theorem describes how an expected value evolves over time.

\begin{theorem}[Ehrenfest's Theorem]
    Let \(\hat{A}\) be a Hermitian operator. Then we have
    \begin{equation}
        \dv{t} \ev{\hat{A}}_\psi = \frac{i}{\hbar} \ev{\comm{\hat{H}}{\hat{A}}}_\psi + \ev{\pdv{\hat{A}}{t}}_\psi.
    \end{equation}
\end{theorem}

\begin{proof}
    We have
    \begin{subequations}
        \begin{align}
            \dv{t} \ev{\hat{A}}_\psi &= \dv{t} \pqty{\psi, \hat{A} \psi} \\
            &= \pqty{\pdv{\psi}{t}, \hat{A} \psi} + \pqty{\psi, \pdv{\hat{A}}{t} \psi} + \pqty{\psi, \hat{A} \pdv{\psi}{t}}.
        \end{align}
    \end{subequations}

    The middle term is precisely \(\ev{\pdv*{\hat{A}}{t}}_\psi\). For the first and last terms, we can use the Schrödinger equation to write
    \begin{equation}
        \pdv{\psi}{t} = - \frac{i}{\hbar} \hat{H} \psi,
    \end{equation}
    and hence
    \begin{subequations}
        \begin{align}
            \pqty{\pdv{\psi}{t}, \hat{A} \psi} &= - \frac{i}{\hbar} \pqty{\hat{H} \psi, \hat{A} \psi} = - \frac{i}{\hbar} \pqty{\psi, \hat{H} \hat{A} \psi} = - \frac{i}{\hbar} \ev{\hat{H} \hat{A}}_\psi, \\
            \pqty{\psi, \hat{A} \pdv{\psi}{t}} &= - \frac{i}{\hbar} \pqty{\psi, \hat{A} \hat{H} \psi} = - \frac{i}{\hbar} \ev{\hat{A} \hat{H}}_\psi.
        \end{align}
    \end{subequations}

    Hence,
    \begin{subequations}
        \begin{align}
            \dv{t} \ev{\hat{A}}_\psi &= - \frac{i}{h} \ev{\hat{H} \hat{A}}_\psi - \frac{i}{\hbar} \ev{\hat{A} \hat{H}}_\psi + \ev{\pdv{\hat{A}}{t}}_\psi \\
            &= \frac{i}{\hbar} \ev{\comm{\hat{H}}{\hat{A}}}_\psi + \ev{\pdv{\hat{A}}{t}}_\psi.
        \end{align}
    \end{subequations}
\end{proof}

\begin{example}[Applications of Ehrenfest's Theorem]
    \begin{itemize}
        \item Take \(\hat{A} = \hat{H}\). Then \(\dv*{E}{t} = 0\) gives the conservation of energy.
        \item Take \(\hat{A} = \hat{p}\). Then
        \begin{equation}
            \dv{t} \ev{\hat{p}}_\psi = - \ev{\dv{U}{x}}_\psi,
        \end{equation}
        which is the quantum analogue of Newton's second law.
        \item Take \(\hat{A} = \hat{x}\). Then
        \begin{equation}
            \dv{t} \ev{\hat{x}}_\psi = \frac{1}{m} \ev{\hat{p}}_\psi,
        \end{equation}
        which is the quantum analogue of the classical relation between position and momentum.
    \end{itemize}
\end{example}